% Cuerpo científico recuperado y distribuido en su residencia terminal.
% Cuerpo editor-ready, sin preámbulo y sin portada autónoma.
% Residencia propuesta: inmediatamente después del teorema del doble círculo
% y antes de Navier--Stokes. No se consume como sustituto de RH ni del
% desarrollo íntegro del doble círculo.

\section{Entrelacement exact de l'holonomie nonadique et de la phase spectrale}
\label{sec:cpe-entrelazamiento-exacto}

Soit \(\mathcal H_W\) l'espace de Hilbert obtenu à partir de la forme positive de
Guinand--Weil, soit \(H_W\) le générateur autoadjoint du groupe de
translations et soit
\[
 \mathcal C_W=(H_W+i/2)(H_W-i/2)^{-1}
\]
sa transformée de Cayley. Pour chaque profondeur \(k\), soit \(Z_k\)
l'ensemble des états cylindriques enrichis, soit \(\mu_k\) sa mesure
projective à support plein et soit
\(\tau_{k+1,k}:Z_{k+1}\to Z_k\) la troncature. On considère
\[
 \mathscr K_k=\ell^2(Z_k)\widehat\otimes\mathcal H_W,
 \qquad
 \Omega_k=\sum_{z\in Z_k}\sqrt{\mu_k(z)}\,e_z,
 \qquad
 J_k\psi=\Omega_k\otimes\psi.
\]
Si \(z'\) est un descendant de \(z\), on prend
\(w_k(z'\mid z)=\sqrt{\mu_{k+1}(z')/\mu_k(z)}\). Pour une route
\(\gamma:z\to z'\), le cocycle entier de mémoire est noté
\(m(\gamma)\), et le transport tordu est défini par
\begin{equation}
 \widetilde U_k(e_z\otimes\psi)
 :=\sum_{\tau_{k+1,k}z'=z}
 w_k(z'\mid z)\,e_{z'}\otimes
 \mathcal C_W^{\,m(z\to z')}\psi.
 \label{eq:cpe-transporte-cilindrico-cayley}
\end{equation}
La projectivité de \(\mu_k\), l'orthogonalité des descendants et
l'unitarité de \(\mathcal C_W\) donnent
\(\widetilde U_k^*\widetilde U_k=I\). L'étiquette complète \(z'\)
conserve le cylindre, le feuillet, l'orientation, la retenue, la route, la frontière et la
multiplicité ; des histoires distinctes ne sont donc pas identifiées lors de l'application de
\eqref{eq:cpe-transporte-cilindrico-cayley}.

\begin{theorem}[Entrelacement de Cayley sous \(1\) tour nonadique complet]
\label{thm:cpe-entrelazamiento-cayley-tpk}
Si \(\Gamma_9\) est l'holonomie nonadique complète, normalisée par
\(m(\Gamma_9)=1\), alors
\begin{equation}
 \boxed{\widetilde U_{\Gamma_9,k}J_k
 =J_{k+9}\mathcal C_W.}
 \label{eq:cpe-entrelazamiento-exacto}
\end{equation}
Pour tout atome spectral \(\psi_\gamma\) de \(H_W\), avec
\(H_W\psi_\gamma=\gamma\psi_\gamma\), on a
\begin{equation}
 \widetilde U_{\Gamma_9^n,k}J_k\psi_\gamma
 =\tau_\gamma^nJ_{k+9n}\psi_\gamma,
 \qquad
 \tau_\gamma=\frac{\gamma+i/2}{\gamma-i/2}\in S^1.
 \label{eq:cpe-orbita-espectral-ledger}
\end{equation}
Dans la publication contractive, la même identité prend la forme
\[
 \widetilde M_{\Gamma_9^n,k}J_k\psi_\gamma
 =\left(\frac59\right)^n\tau_\gamma^n
 J_{k+9n}\psi_\gamma.
\]
\end{theorem}

\begin{proof}
Le tour \(\Gamma_9\) revient à la phase visible mais augmente d'une
unité la mémoire, de sorte que \(m(\Gamma_9)=1\). En sommant tous les
descendants de \(\Omega_k\), les poids se télescopent avec la mesure projective
et produisent \(\Omega_{k+9}\), tandis que la fibre de Weil reçoit exactement
une puissance de \(\mathcal C_W\). Cela prouve
\(\widetilde U_{\Gamma_9,k}J_k=J_{k+9}\mathcal C_W\). L'itération utilise
l'additivité \(m(\Gamma_9^n)=n\). La dernière formule découle de
\(\widetilde M_{\Gamma_9,k}=(5/9)\widetilde U_{\Gamma_9,k}\).
\end{proof}

Le théorème n'identifie pas un zéro à une phase isolée du calendrier. Un
tour complet applique une fois le caractère de Cayley ; chaque zéro détermine
la loi angulaire suivant laquelle une histoire traverse tous les tours. En
particulier, le retour de phase et le retour d'état demeurent distincts :
\[
 g(K+9)=g(K),
 \qquad B_{K+9}\ne B_K\quad\text{en général}.
\]

\section{Convergence en norme des matrices spectrales finies}
\label{sec:cpe-convergencia-resolvente}

Soit
\[
 R_W=(H_W-iI)^{-1}.
\]
Le spectre de \(H_W\) est discret et \(|\gamma|\to\infty\), donc
\(R_W\) est compact. On prend une famille dénombrable
\(\{f_j\}_{j\ge1}\subset C_c^\infty(\mathbb R)\) à paramètres rationnels,
dense dans le noyau de tests et choisie sans consulter de hauteurs spectrales.
Après passage au quotient par le noyau de la forme de Weil et orthonormalisation, on
obtient des sous-espaces
\[
 \mathcal V_N=\operatorname{span}\{[f_1],\ldots,[f_N]\},
 \qquad P_N:\mathcal H_W\to\mathcal V_N,
\]
avec \(P_N\to I\) fortement. Définissons
\begin{equation}
 R_N=P_NR_WP_N.
 \label{eq:cpe-resolvente-finito}
\end{equation}

\begin{theorem}[Approximation compacte sans introduire les zéros]
\label{thm:cpe-convergencia-norma}
La suite \(R_N\) converge en norme vers \(R_W\) :
\begin{equation}
 \boxed{\lim_{N\to\infty}\|R_W-R_N\|=0.}
 \label{eq:cpe-convergencia-norma}
\end{equation}
Par conséquent, chaque valeur propre non nulle
\(\lambda_\gamma=(\gamma-i)^{-1}\) de \(R_W\) est limite, avec sa
multiplicité algébrique, de valeurs propres des matrices finies de
\(R_N\). Les hauteurs sont retrouvées au moyen de
\begin{equation}
 \gamma=i+\lambda_\gamma^{-1}.
 \label{eq:cpe-recuperacion-altura}
\end{equation}
\end{theorem}

\begin{proof}
Si \(K\) est compact et \(P_N\to I\) fortement, la convergence est uniforme
sur le compact \(K(B_1)\), où \(B_1\) est la boule unité. Par conséquent,
\[
 \|(I-P_N)K\|\longrightarrow0.
\]
En appliquant le même argument à \(K^*\), on obtient
\(\|K(I-P_N)\|\to0\). Ainsi,
\[
 \|K-P_NKP_N\|
 \leq\|(I-P_N)K\|+\|P_NK(I-P_N)\|\longrightarrow0.
\]
En prenant \(K=R_W\), on prouve \eqref{eq:cpe-convergencia-norma}. La
stabilité des valeurs propres isolées d'opérateurs compacts sous des
perturbations en norme et le calcul fonctionnel de Riesz conservent leur
multiplicité. L'identité \eqref{eq:cpe-recuperacion-altura} est l'inverse
de l'application spectrale de la résolvante.
\end{proof}

Les entrées de \(R_N\) sont calculées à partir de la forme premier--gamma--polaire :
\[
 \langle R_W[f_j],[f_k]\rangle_W
 =i\int_0^\infty e^{-t}
 \mathscr I_{\rm GW}(T_{-t}f_j,f_k)\,dt.
\]
Le théorème n'utilise donc pas de table de zéros pour former les matrices.

\subsection{Certification a posteriori des intervalles et des multiplicités}

Soit \((\lambda_N,v_N)\) un vecteur propre normalisé de \(R_N\) et soit
\[
 \varepsilon_N(v_N)
 :=\|(I-P_N)R_Wv_N\|.
\]
Comme \(R_W\) est normal,
\begin{equation}
 \operatorname{dist}(\lambda_N,\sigma(R_W))
 \leq\varepsilon_N(v_N).
 \label{eq:cpe-residuo-inclusion}
\end{equation}
Les canaux premier, gamma et polaire fournissent des bornes par intervalles séparées
pour \(\varepsilon_N\) ; leurs queues sont sommées avant d'émettre l'intervalle.
De plus, si \(\Gamma\) est un contour fermé contenu dans l'ensemble résolvant de
\(R_N\) et
\begin{equation}
 \sup_{z\in\Gamma}
 \|(z-R_N)^{-1}\|\,\|R_W-R_N\|<1,
 \label{eq:cpe-criterio-riesz}
\end{equation}
alors les projections de Riesz de \(R_N\) et de \(R_W\) délimitées par
\(\Gamma\) ont le même rang. Les équations
\eqref{eq:cpe-residuo-inclusion} et \eqref{eq:cpe-criterio-riesz} certifient,
respectivement, l'inclusion et la complétude des atomes contenus dans chaque
fenêtre. L'arithmétique des intervalles consolide les décimales ; elle n'ajoute pas d'application
mathématique absente.

\section{Trace couplée du spectre de Weil et de la graduation Moonshine}
\label{sec:cpe-traza-rh-moonshine}

La trace produit
\[
 \operatorname{Tr}_{\mathcal H_W}
 \bigl(e^{-aH_W^2}\mathcal C_W^n\bigr)
 \operatorname{Tr}_{V^\natural}
 \bigl(g e^{-\beta(L_0-1)}\bigr)
\]
superpose deux observables, mais ne les couple pas. La mémoire du TPK et la
graduation de \(V^\natural\) permettent de construire le couplage sans postuler
d'action du Monstre sur des chiffres ni sur des hauteurs.

Une fois fixée la jauge exceptionnelle \(\chi_{\rm exc}\), on oriente conjointement
l'origine de la mémoire et l'origine de la graduation. Soit
\(\mathcal H_{\rm mem}=\ell^2(\mathbb N_0)\), avec
\(N_{\rm mem}|n\rangle=n|n\rangle\), et soit
\[
 V^\natural=\widehat\bigoplus_{r\ge0}V^\natural_r,
 \qquad L_0|_{V^\natural_r}=rI.
\]
Notons \(P_r^\natural\) la projection sur
\(V^\natural_r\). Le projecteur d'égalité des degrés est
\begin{equation}
 \Pi_{\chi_{\rm exc}}
 :=\sum_{n\ge0}|n\rangle\langle n|
 \otimes I_{\mathcal H_W}\otimes P_n^\natural.
 \label{eq:cpe-proyector-diagonal-grados}
\end{equation}
La somme converge fortement et définit une projection orthogonale. Sur le
produit, on définit le caractère de mémoire
\begin{equation}
 \mathcal C_W^{N_{\rm mem}}
 :=\sum_{n\ge0}|n\rangle\langle n|
 \otimes\mathcal C_W^n.
 \label{eq:cpe-cayley-memoria}
\end{equation}

\begin{definition}[Trace diagonale RH--Moonshine]
\label{def:cpe-traza-diagonal}
Pour \(a,b,\beta>0\) et \(g\in\mathbb M\), on définit
\begin{equation}
\begin{aligned}
 \Theta^{\Delta}_{g,\chi_{\rm exc}}(a,b,\beta)
 :={}&\operatorname{Tr}\!\left[
 \Pi_{\chi_{\rm exc}}
 \Bigl(
 e^{-bN_{\rm mem}}
 \mathcal C_W^{N_{\rm mem}}e^{-aH_W^2}
 \otimes g e^{-\beta(L_0-1)}
 \Bigr)
 \Pi_{\chi_{\rm exc}}
 \right].
 \label{eq:cpe-traza-diagonal-def}
\end{aligned}
\end{equation}
\end{definition}

\begin{theorem}[Existence et développement non factorisé]
\label{thm:cpe-traza-diagonal}
La trace \eqref{eq:cpe-traza-diagonal-def} est absolument convergente et
satisfait
\begin{equation}
 \boxed{\Theta^{\Delta}_{g,\chi_{\rm exc}}(a,b,\beta)
 =\sum_{n\ge0}e^{-bn}e^{-\beta(n-1)}
 \operatorname{Tr}_{\mathcal H_W}
 \bigl(e^{-aH_W^2}\mathcal C_W^n\bigr)
 \operatorname{Tr}_{V^\natural_n}(g).}
 \label{eq:cpe-traza-diagonal-expansion}
\end{equation}
Après la reconnaissance spectrale,
\begin{equation}
 \Theta^{\Delta}_{g,\chi_{\rm exc}}(a,b,\beta)
 =\sum_{n\ge0}\sum_\gamma
 m_\gamma e^{-a\gamma^2}
 e^{-bn-\beta(n-1)}\tau_\gamma^n
 \operatorname{Tr}_{V^\natural_n}(g).
 \label{eq:cpe-traza-diagonal-espectral}
\end{equation}
Cette expression n'est pas le produit d'une somme spectrale et d'une série de
McKay--Thompson : le même indice \(n\) sélectionne simultanément la puissance
du caractère de mémoire et le degré monstrueux.
\end{theorem}

\begin{proof}
Pour \(a>0\), la loi de comptage spectral implique
\(\sum_\gamma m_\gamma e^{-a\gamma^2}<\infty\). Chaque espace
\(V^\natural_n\) est de dimension finie et les coefficients des caractères
de Moonshine croissent sous-exponentiellement en \(n\). Le facteur
\(e^{-(b+\beta)n}\) assure la convergence absolue de la série conjointe.
Le développement est obtenu en insérant les résolutions spectrales de
\(N_{\rm mem}\), \(H_W\) et \(L_0\) et en utilisant
\eqref{eq:cpe-proyector-diagonal-grados}. La présence du projecteur diagonal
élimine les termes de degrés différents et produit une seule somme en
\(n\), au lieu des deux sommes indépendantes de la trace produit.
\end{proof}

Le couplage dépend de la jauge \(\chi_{\rm exc}\), qui fixe l'origine,
l'orientation et l'alignement de la lecture exceptionnelle. Un choix conjugué
produit une trace conjuguée ou une réindexation déclarée ; il ne modifie pas le
théorème de Riemann--Weil ni l'entrelacement
\eqref{eq:cpe-entrelazamiento-exacto}.

